% 光滑函数（综述）
% license CCBYSA3
% type Wiki

本文根据 CC-BY-SA 协议转载翻译自维基百科\href{https://en.wikipedia.org/wiki/Smoothness}{相关文章}

在数学分析中，函数的光滑性是一个由其定义域上连续导数的阶数（可微类，differentiability class）来衡量的性质。[1]

一个属于$C^k$类的函数，意味着它至少是$k$阶光滑的；换句话说，$C^k$类函数是在其定义域上具有连续的$k$阶导数的函数。

一个属于$C^\infty$类的函数（也称为$C^\infty$函数，读作“C-infinity function”），是一个无穷可微函数，即拥有所有阶数导数的函数（这也意味着所有这些导数都是连续的）。

通常情况下，术语光滑函数指的就是一个$C^\infty$函数。但在某些问题的语境中，它也可能表示“具有足够阶可微性”的函数。
\subsection{可微类}
可微类是根据函数导数的性质对函数进行分类的一种方式，它衡量的是函数存在并且连续的最高阶导数的阶数。

考虑实数轴上的一个开集$U$，以及一个定义在$U$上、取实值的函数$f$。设$k$是一个非负整数。如果函数$f$的各阶导数$f', f'', \dots, f^{(k)}$都存在并且在$U$上连续，则称函数$f$属于$C^k$类。如果 $f$ 在 $U$ 上是 $k$ 阶可微的，那么它至少属于$C^{k-1}$类，因为它的 $f', f'', \dots, f^{(k-1)}$都在$U$上连续。如果函数 $f$ 在 $U$ 上有任意阶导数，则称它是无穷可微的、光滑的，或者是$C^\infty$ 类。这意味着它的所有导数在$U$上都是连续函数。[2]如果函数$f$是光滑的（即 $f \in C^\infty$），并且它在定义域内任意一点的 Taylor 级数在该点的某个邻域内收敛并等于函数本身，则称$f$属于$C^\omega$类，也称为解析函数。存在一些函数是光滑但不是解析的，因此 $C^\omega$ 严格包含在 $C^\infty$ 之内。凸起函数就是这种性质的例子。

换句话说：$C^0$类是所有连续函数的集合；$C^1$类是所有一阶导数存在且连续的函数，这类函数称为连续可微函数。因此，$C^1$函数正好是导数存在且导数属于$C^0$类的函数。一般来说，可递归地定义这些类：$C^0$类：所有连续函数的集合；$C^k$类（$k>0$）：所有可微函数的集合，其导数属于$C^{k-1}$ 类。由此可以得出：对于每个$k>0$，有$C^k \subset C^{k-1},$且严格包含（$C^k \subsetneq C^{k-1}$）。无穷可微函数的类$C^\infty$是所有非负整数$k$的$C^k$类的交集。
